% !TeX program = xelatex
\documentclass[11pt,a4paper,oneside]{report}
\usepackage{fontspec}
\setmainfont{Liberation Serif}
\setsansfont{Liberation Sans}
\setmonofont{Liberation Mono}[Scale=MatchLowercase]
\usepackage[vietnamese]{babel}
\usepackage[a4paper,left=2.7cm,right=2.3cm,top=2.35cm,bottom=2.4cm,headheight=15pt]{geometry}
\usepackage{amsmath,amssymb,graphicx,calc,booktabs,longtable,tabularx,array,multirow}
\usepackage{float,pdflscape,xcolor,enumitem,indentfirst,fancyhdr,titlesec,caption}
\usepackage{tikz}
\usetikzlibrary{calc,shapes.geometric,arrows.meta,positioning,shadows}
\usepackage{fvextra}
\usepackage{hyperref,xurl,bookmark}
\graphicspath{{../assets/}{../assets/screenshots/}{../../diagrams/}{./}}
\definecolor{hubblue}{HTML}{124159}
\definecolor{hublink}{HTML}{165D75}
\hypersetup{unicode=true,colorlinks=true,allcolors=hublink,bookmarksnumbered=true,
  pdftitle={Smart E-Mobility Hub - Báo cáo tổng kết Submission 4},
  pdfauthor={Nhóm BEE BRAVE, L02, CO3001, HK261}}
\setcounter{tocdepth}{1}
\setcounter{secnumdepth}{2}
\setlength{\emergencystretch}{3em}
\newlength{\reporttablewidth}
\pdfstringdefDisableCommands{\def\allowbreak{}}
\captionsetup{font=small,labelfont=bf}
\titleformat{\chapter}[hang]{\normalfont\LARGE\bfseries\color{hubblue}}{\thechapter.}{0.7em}{}
\titlespacing*{\chapter}{0pt}{0pt}{16pt}
\fancypagestyle{plain}{%
  \fancyhf{}%
  \fancyhead[L]{\small CO3001 -- Công nghệ Phần mềm -- HK261}%
  \fancyhead[R]{\small BEE BRAVE | Submission 4}%
  \fancyfoot[C]{\thepage}%
}
\pagestyle{plain}
\providecommand{\tightlist}{\setlength{\itemsep}{0pt}\setlength{\parskip}{0pt}}
\providecommand{\pandocbounded}[1]{#1}
% Pandoc images share a maximum size and retain their aspect ratio.
\setkeys{Gin}{width=\linewidth,height=.71\textheight,keepaspectratio}
\DefineVerbatimEnvironment{Highlighting}{Verbatim}{breaklines,fontsize=\small}
\newcommand{\evidencefigure}[3]{%
  \begin{figure}[H]\centering
  \includegraphics[width=\linewidth,height=.71\textheight,keepaspectratio]{#1}
  \caption{#2}\label{#3}\end{figure}}
\newcommand{\demoscreen}[4]{%
  \clearpage\section{#1}#2
  \evidencefigure{figures/demo/#3}{#4}{fig:demo-#3}}
\begin{document}
\shorthandoff{"}
\pagenumbering{gobble}
\begin{titlepage}
  \begin{tikzpicture}[remember picture,overlay]
    \draw[hubblue,line width=1.7pt]
      ($(current page.north west)+(1.35cm,-1.35cm)$) rectangle
      ($(current page.south east)+(-1.35cm,1.35cm)$);
  \end{tikzpicture}
  \centering
  {\large\textbf{ĐẠI HỌC QUỐC GIA THÀNH PHỐ HỒ CHÍ MINH}}\\[0.28cm]
  {\Large\textbf{TRƯỜNG ĐẠI HỌC BÁCH KHOA}}\\[0.18cm]
  {\large\textbf{KHOA KHOA HỌC VÀ KỸ THUẬT MÁY TÍNH}}\\[0.8cm]
  \includegraphics[width=.24\textwidth,height=3.6cm,keepaspectratio]{../assets/mau.png}\\[0.75cm]
  \rule{.86\textwidth}{.45pt}\\[0.85cm]
  {\LARGE\bfseries BÁO CÁO TỔNG KẾT BÀI TẬP LỚN}\\[0.3cm]
  {\Large\bfseries SUBMISSION \#4}\\[0.8cm]
  {\large\bfseries CÔNG NGHỆ PHẦN MỀM (CO3001) -- HK261}\\[0.65cm]
  {\huge\bfseries\color{hubblue} SMART E-MOBILITY HUB}\\[0.25cm]
  {\large Hệ thống điều phối phương tiện điện\\
    trong Khu đô thị ĐHQG-HCM}\\[0.75cm]
  \rule{.86\textwidth}{.45pt}\\[0.7cm]
  \begin{tabular}{@{}ll@{}}
    \textbf{Giảng viên hướng dẫn} & TS. Trần Thị Ngọc Trâm \\
    \textbf{Lớp học phần} & L02 \\
    \textbf{Nhóm thực hiện} & BEE BRAVE (7 thành viên)
  \end{tabular}
  \vfill
  {\large Thành phố Hồ Chí Minh, tháng 11 năm 2026}
\end{titlepage}

\pagenumbering{roman}
\chapter*{Danh sách thành viên và phân công}
\addcontentsline{toc}{chapter}{Danh sách thành viên và phân công}
\noindent\begin{tabularx}{\linewidth}{@{}ll l X@{}}
\toprule
\textbf{Phân hệ} & \textbf{Thành viên} & \textbf{MSSV} & \textbf{Phạm vi phụ trách} \\
\midrule
SS1 & Nguyễn Thành Đạt & 2410709 & Đặt và quản lý xe dùng chung, trưởng nhóm \\
SS2 & Nguyễn Anh Tài & 2413035 & Xe cá nhân, đặt chỗ đỗ, check-in và check-out \\
SS3 & Phạm Tiến Đạt & 2410724 & Giám sát mạng lưới Hub \\
SS4 & Đặng Hoàng Quý Nhân & 2412401 & Điều phối phương tiện và xử lý sự cố \\
SS5 & Dương Đăng Khoa & 2411589 & Lập lịch và quản lý phiên sạc \\
SS6 & Hoàng Văn Tấn & 2213074 & Mô phỏng nhu cầu và dự báo \\
SS7 & Nguyễn Trung Nguyên & 2011710 & Mô phỏng sự cố hạ tầng và điều hướng \\
\bottomrule
\end{tabularx}

\tableofcontents
\listoffigures
\clearpage
\pagenumbering{arabic}
\chapter*{Tóm tắt}
\addcontentsline{toc}{chapter}{Tóm tắt}
Smart E-Mobility Hub quản lý phương tiện điện, chỗ đỗ và cổng sạc tại sáu Hub mô phỏng
trong Khu đô thị ĐHQG-HCM. Sinh viên đặt và trả xe dùng chung hoặc gửi và sạc xe cá nhân.
Nhân viên vận hành giám sát trạng thái, điều phối xe và kiểm tra các phương án What-if.
Bản MVP dùng Python và Streamlit. Bộ phát IoT độc lập giúp kiểm tra cách ứng dụng
phản ứng khi năng lực sạc thay đổi trong lúc người dùng đang thao tác.
\par
Báo cáo hợp nhất nội dung của ba lần nộp trước trong một bản LaTeX.
Phần I trình bày yêu cầu và use case. Phần II trình bày mockup, sơ đồ tuần tự,
hoạt động và trạng thái. Phần III mô tả triển khai, cấu trúc module, lớp,
phương thức. Phần IV ghi lại quy trình demo và khai báo sử dụng AI.
\par
Các sơ đồ thiết kế có thể chứa service hoặc giao tiếp được đề xuất trong Submission 1--3.
Khi cần xác định hành vi của bản trình diễn, đối chiếu chương~\ref{chap:runtime}
và mã nguồn trong gói nộp. Phần demo dùng ảnh chụp trình duyệt từ
lần chạy ngày 24/11/2026. Phạm vi chứng minh là chạy local.

\part{Phân tích và đặc tả yêu cầu}
\chapter{Đặc tả hệ thống}
\input{build/generated/requirements}
\chapter{Mô hình use case toàn hệ thống}
\label{chap:usecases}
\evidencefigure{../../diagrams/system/smartEhub.png}{Use case tổng thể của bảy phân hệ.}{fig:usecases}
\evidencefigure{../../diagrams/telemetry/usecase.png}{Bộ phát cảm biến và nhân viên vận hành trong luồng IoT mô phỏng.}{fig:iot-actors}
\section{Danh mục use case}
Mã định danh lấy từ ma trận truy vết chung. Bản tin cảm biến thuộc luồng cảnh báo của
\texttt{UC\_SS3\_02}.
\input{build/generated/use-case-index}
\section{Kịch bản use case theo phân hệ}
Các bảng dưới đây tóm lược đặc tả cá nhân trong \path{docs/diagrams/Use-case detail-scenario/}.
Mã và thứ tự use case theo ma trận chung. Xác thực, barie và thiết bị được mô phỏng
trong các thao tác của MVP. Luồng nghiệp vụ chạy trong từng phiên Streamlit.

\newcommand{\casebrief}[7]{%
  \subsection[#1: #2]{#1: #2}%
  \noindent\begin{tabularx}{\linewidth}{@{}>{\bfseries\raggedright\arraybackslash}p{.22\linewidth}X@{}}
  \toprule
  Phụ trách & #3 \\
  Tác nhân & #4 \\
  Điều kiện trước & #5 \\
  Luồng chính & #6 \\
  \midrule
  Điều kiện sau và nhánh xử lý & #7 \\
  \bottomrule
  \end{tabularx}\par}

\casebrief{UC\_SS1\_01}{Tra cứu, đặt và nhận xe}
{Nguyễn Thành Đạt}{Sinh viên, ví nội bộ và khóa mô phỏng.}
{Hub có xe dùng chung khả dụng đủ pin, sinh viên không có lượt đang hoạt động và ví đủ cọc.}
{Chọn Hub và xe, nhập voucher nếu có, xác nhận giữ xe. Hệ thống khóa cọc, giữ xe 30 phút và cấp PIN. Sinh viên xác nhận nhận xe bằng PIN.}
{Xe sang \texttt{IN\_USE}, chỗ tại Hub xuất phát được giải phóng. PIN sai giữ nguyên lượt. Tại đúng hạn mà chưa nhận, lượt hết hạn và hệ thống quyết toán no-show.}

\casebrief{UC\_SS1\_02}{Trả xe và quyết toán}
{Nguyễn Thành Đạt}{Sinh viên, ví nội bộ và Hub đích.}
{Lượt đang di chuyển, Hub trả có chỗ còn trống hoặc có chỗ đã giữ cho lượt đó.}
{Chọn Hub trả và xác nhận. Hệ thống cập nhật vị trí xe, tình trạng pin, thời lượng và quãng đường mô phỏng, sau đó quyết toán cước và tiền cọc.}
{Lượt hoàn tất và lưu hóa đơn. Nếu Hub đầy, chọn Hub khác hoặc nhận gợi ý. Nếu thiếu tiền, ghi nợ phần còn lại. Lượt đã hoàn tất không được quyết toán lại.}

\casebrief{UC\_SS1\_03}{Hủy giữ xe}
{Nguyễn Thành Đạt}{Sinh viên và ví nội bộ.}
{Lượt giữ xe còn ở trạng thái \texttt{RESERVED}.}
{Sinh viên xác nhận hủy trước hạn. Hệ thống trả xe về danh sách khả dụng, giải phóng chỗ đã giữ ở đích và trả cọc.}
{Lượt chuyển \texttt{CANCELLED}. Từ mốc hết hạn, áp dụng no-show. Lượt đã nhận xe không thể hủy theo quy trình này.}

\casebrief{UC\_SS2\_01}{Đăng ký xe điện cá nhân}
{Nguyễn Anh Tài}{Sinh viên.}
{Có phiên người dùng demo để liên kết hồ sơ.}
{Nhập biển số, loại xe và thông tin pin. Hệ thống chuẩn hóa biển số, kiểm tra hợp lệ và cấp mã xe.}
{Hồ sơ được lưu và có thể chọn khi đặt dịch vụ. Hãng xe có thể bỏ trống. Biển số trùng hoặc loại xe/dung lượng pin không hợp lệ bị từ chối.}

\casebrief{UC\_SS2\_02}{Đặt chỗ đỗ và nhu cầu sạc}
{Nguyễn Anh Tài}{Sinh viên và ví nội bộ.}
{Xe thuộc người đặt, Hub có chỗ, thời điểm đến hợp lệ và ví đủ cọc.}
{Chọn xe, Hub, thời điểm đến và nhu cầu sạc. Hệ thống giữ chỗ đỗ, khóa cọc và lập lượt đặt với hạn check-in 30 phút tính từ thời điểm đến.}
{Lượt ở trạng thái \texttt{RESERVED}. Chỉ đỗ thì bỏ nhu cầu sạc. Hủy trước hạn được hoàn cọc. Hết hạn thì giải phóng chỗ và xử lý no-show.}

\casebrief{UC\_SS2\_03}{Check-in xe cá nhân}
{Nguyễn Anh Tài}{Sinh viên, cổng vào mô phỏng và bộ lập lịch SS5.}
{Có lượt giữ cho đúng xe và Hub trong cửa sổ check-in hợp lệ.}
{Xác nhận check-in. Hệ thống ghi giờ vào, chuyển chỗ giữ thành chỗ có xe. Nếu có nhu cầu sạc thì đưa yêu cầu vào SS5 để cấp cổng khả dụng.}
{Lượt chuyển \texttt{CHECKED\_IN}. Xe đợi nếu chưa cấp được cổng. Đến sớm bị từ chối và đến đúng mốc hết hạn bị xử lý no-show.}

\casebrief{UC\_SS2\_04}{Check-out và thanh toán tổng hợp}
{Nguyễn Anh Tài}{Sinh viên, cổng ra mô phỏng, ví nội bộ và SS5.}
{Xe có lượt \texttt{CHECKED\_IN} đang hoạt động.}
{Xác nhận rời trạm. Hệ thống chốt thời gian đỗ và phí theo giờ bắt đầu, kết thúc phiên sạc nếu có, tổng hợp điện năng, phí quá giờ và quyết toán ví.}
{Chỗ đỗ và cổng đã cấp được giải phóng, hóa đơn được lưu, lượt hoàn tất. Chỉ đỗ thì không thu phí sạc. Thiếu số dư thì ghi nợ. Check-out lặp bị từ chối.}

\casebrief{UC\_SS3\_01}{Giám sát mạng lưới Hub}
{Phạm Tiến Đạt}{Nhân viên vận hành, quản trị viên theo thiết kế và bộ phát cảm biến mô phỏng.}
{Dữ liệu Hub đã được nạp. Bật IoT nếu cần nhận bản tin trên giao diện.}
{Mở dashboard, theo dõi sức chứa, xe, cổng sạc, doanh thu, giao dịch và công suất. Khi bật IoT, mỗi phiên định kỳ đọc feed và làm mới màn hình đang mở.}
{Hiển thị snapshot của phiên. Nếu chưa bật feed, màn hình dùng dữ liệu mô phỏng hiện có. Dữ liệu ví và phiên sạc không được chia sẻ giữa các cửa sổ.}

\casebrief{UC\_SS3\_02}{Cảnh báo chiếm dụng Hub và sự cố IoT}
{Phạm Tiến Đạt}{Nhân viên vận hành, kỹ thuật viên và bộ phát cảm biến mô phỏng.}
{Có dữ liệu sức chứa, ngưỡng cảnh báo và trạng thái cổng của Hub.}
{Dashboard phân mức lấp đầy ở 70, 85 và 100 phần trăm. Bản tin hợp lệ cập nhật trạng thái cổng. Operator kiểm tra cảnh báo và chuyển sang màn hình xử lý phù hợp.}
{Hub đầy không nhận đặt chỗ mới theo điều kiện sức chứa. Bản tin trùng ID, cũ hoặc sai đích không tạo thêm tác động vận hành.}

\casebrief{UC\_SS4\_01}{Điều phối phương tiện giữa các Hub}
{Đặng Hoàng Quý Nhân}{Nhân viên vận hành và hệ thống quản lý Hub.}
{Hub nguồn có xe chung khả dụng. Hub đích khác nguồn và có đủ chỗ giữ.}
{Chọn Hub nguồn, Hub đích và xe. Tạo lệnh giữ chỗ ở đích và đưa xe vào trạng thái trung chuyển. Sau khi xe đến, xác nhận tiếp nhận.}
{Xe được ghi nhận tại Hub đích, chỗ giữ được chuyển thành chỗ có xe. Thiếu xe hợp lệ hoặc thiếu chỗ thì không lập lệnh.}

\casebrief{UC\_SS4\_02}{Báo cáo và xử lý sự cố phương tiện}
{Đặng Hoàng Quý Nhân}{Nhân viên vận hành và kỹ thuật viên.}
{Phương tiện cần xử lý được nhận diện tại Hub.}
{Nhập đối tượng, mức độ và mô tả lỗi. Hệ thống ghi nhận và cách ly xe. Kỹ thuật viên tiếp nhận rồi xác nhận đã khắc phục.}
{Nhật ký ghi các bước xử lý và tài nguyên được khôi phục theo thao tác xác nhận. Sự cố chưa xử lý vẫn hiển thị trong danh sách theo dõi.}

\casebrief{UC\_SS5\_01}{Ưu tiên sạc và cấp cổng}
{Dương Đăng Khoa}{Nhân viên vận hành sạc, SS2 và bộ lập lịch.}
{Có yêu cầu chờ và dữ liệu tài nguyên, điểm pin và công suất Hub.}
{Tính điểm theo mức pin và ưu tiên booking, dùng thời điểm yêu cầu để phá hòa. Cấp cổng theo thứ tự ưu tiên trong giới hạn số cổng và công suất.}
{Xe được cấp cổng có phiên sạc. Xe còn lại tiếp tục chờ. Nếu mất cổng, phiên bị ảnh hưởng được xử lý và giữ phần điện, phí đã tích lũy.}

\casebrief{UC\_SS5\_02}{Theo dõi sạc và phí chiếm dụng}
{Dương Đăng Khoa}{Nhân viên vận hành sạc và người sử dụng dịch vụ.}
{Có phiên sạc hoặc yêu cầu chờ tại Hub.}
{Theo dõi phiên, dùng bước mô phỏng để cập nhật pin và điện năng. Khi đạt mức pin mục tiêu thì bắt đầu ân hạn rút sạc. Xác nhận rút sạc để trả cổng.}
{Phiên lập biên nhận hoặc chuyển số liệu vào hóa đơn SS2. Nếu giữ cổng sau 15 phút ân hạn thì phát sinh phí. Sự cố cổng không làm mất số điện đã nạp.}

\casebrief{UC\_SS6\_01}{Mô phỏng đợt tăng nhu cầu Metro}
{Hoàng Văn Tấn}{Nhà phân tích hoặc nhân viên vận hành, bộ mô phỏng.}
{Đã nạp snapshot Hub và có số lượng sinh viên giả lập.}
{Chọn Hub và quy mô đợt đến. Sinh preview trên bản sao, so sánh số người tiếp nhận được và số chưa phục vụ, sau đó quyết định áp dụng.}
{Preview chưa đổi trạng thái vận hành. Có thể đổi Hub hoặc quy mô. Apply bị từ chối nếu phiên bản dữ liệu đã đổi.}

\casebrief{UC\_SS6\_02}{Dự báo nhu cầu và hàng đợi 24 giờ}
{Hoàng Văn Tấn}{Nhà phân tích hoặc người hoạch định vận hành.}
{Có bộ tham số nhu cầu, công suất phục vụ giả định và khoảng dự báo.}
{Điều chỉnh tỉ lệ đến và hệ số tăng. Bộ dự báo sinh chuỗi theo giờ và hiển thị lượng đến, hàng chờ, thời gian chờ và doanh thu giả định.}
{Kết quả chỉ phục vụ phân tích. Tăng nhu cầu để xem kịch bản tải cao hoặc giảm để so sánh. Dự báo không làm thay đổi ví hay doanh thu đã thu.}

\casebrief{UC\_SS7\_01}{Mô phỏng lỗi cổng sạc}
{Nguyễn Trung Nguyên}{Kỹ thuật viên hoặc nhân viên vận hành, bộ mô phỏng hạ tầng.}
{Có snapshot Hub và số cổng với trạng thái đã biết.}
{Chọn Hub, số cổng hỏng rồi tạo preview. Xem tác động tới cổng và phiên sạc, xem gợi ý Hub khác và xác nhận Apply nếu phù hợp.}
{Apply cập nhật trạng thái phiên khi dữ liệu khớp phiên bản gốc. Nếu có bản tin làm phiên bản đổi, chạy lại preview. Có thể khôi phục cổng qua thao tác mô phỏng.}

\casebrief{UC\_SS7\_02}{Gợi ý Hub thay thế và điều hướng}
{Nguyễn Trung Nguyên}{Sinh viên, nhân viên vận hành và bộ định tuyến.}
{Hub ban đầu thiếu chỗ hoặc năng lực sạc cần thiết, có dữ liệu trạng thái và vị trí Hub lân cận.}
{Từ kết quả cảnh báo hoặc preview, lọc và xếp các Hub còn tài nguyên theo khoảng cách. Khi sinh viên xác nhận chuyển Hub cho lượt thuê, kiểm tra lại chỗ trước khi chuyển giữ.}
{Thay đổi đích chỉ xảy ra khi xác nhận hợp lệ. Có thể chọn Hub khác thủ công. Nếu phương án hết chỗ hoặc không có ứng viên, giữ nguyên lượt hiện có và chọn lại.}


\part{Thiết kế giao diện và tương tác}
\input{build/generated/interaction}
\chapter{Tương tác với bộ phát cảm biến}
\label{chap:iot-interaction}
Bản tin cảm biến đi từ tiến trình phát độc lập đến feed SQLite. Mỗi phiên giao diện
đọc feed theo chu kỳ 2 giây khi bật IoT. Mã bản tin trùng bị loại ở đầu ghi.
Khi đọc, hệ thống kiểm tra Hub, chỉ số cổng và thứ tự thời gian. Chỉ bản tin hợp lệ
mới cập nhật snapshot và tăng phiên bản trạng thái. SS3 hiển thị cảnh báo và SS7
dùng phiên bản mới để kiểm tra preview.
\evidencefigure{../../diagrams/telemetry/flow.png}{Sequence xử lý bản tin IoT mô phỏng.}{fig:iot-sequence}
\evidencefigure{../../diagrams/telemetry/activity.png}{Activity tiếp nhận và áp dụng bản tin.}{fig:iot-activity}
\part{Thiết kế kiến trúc và hiện thực}
\input{build/generated/architecture}
\chapter{Hiện thực tại revision bàn giao}
\label{chap:runtime}
\section{Kiến trúc chạy local}
Trình duyệt truy cập ứng dụng Streamlit. \path{src/app.py} điều hướng sang các màn hình
trong \path{src/components/}. Các module \path{src/core/} xử lý nghiệp vụ. Seed sáu Hub
nằm trong \path{src/data/mock_hubs.json}. Ví, lượt đặt, phương tiện và phiên sạc được
lưu trong trạng thái phiên. Bộ phát \path{scripts/iot_producer.py} ghi bản tin vào SQLite.
Khi bật IoT, mỗi phiên đọc và áp dụng feed lên snapshot của mình.
Sơ đồ triển khai khi bật IoT được trình bày ở Hình~\ref{fig:s3-1}.
\evidencefigure{../../diagrams/telemetry/classes.png}{Hợp đồng và mô hình dữ liệu của luồng IoT.}{fig:runtime-classes}
\begin{center}
\begin{tabularx}{\linewidth}{@{}p{.25\linewidth}X@{}}
\toprule
\textbf{Module} & \textbf{Trách nhiệm chính}\\
\midrule
\path{core/booking.py} & Giữ xe dùng chung, nhận và trả, hủy hoặc tự hết hạn, chuyển chỗ giữ và tạo hóa đơn.\\
\path{core/personal_ev.py} & Đăng ký xe cá nhân, giữ chỗ đỗ, check-in, gửi yêu cầu sạc và check-out.\\
\path{core/scheduler.py} & Xếp hàng sạc, cấp cổng trong giới hạn công suất, tích lũy điện năng và giải phóng cổng.\\
\path{core/billing.py} & Khóa và giải phóng cọc, ghi sổ cái, quyết toán và ghi nợ.\\
\path{core/telemetry.py} & Lưu, kiểm tra và đọc lại feed. Áp dụng thay đổi trước khi hiển thị phân hệ đang mở.\\
\path{core/simulator.py} & Dự báo nhu cầu và tạo snapshot What-if. Preview chứa phiên bản dữ liệu để kiểm tra khi Apply.\\
\bottomrule
\end{tabularx}
\end{center}

\section{Quy tắc đầu cuối của bản demo}
\begin{itemize}
\item \textbf{SS1 và SS2.} Cọc giữ chỗ là 50.000~đ. SS1 giữ xe trong 30 phút
kể từ lúc đặt. SS2 chấp nhận check-in trong 30 phút tính từ giờ hẹn.
No-show thu 10.000~đ và hoàn 40.000~đ. Hóa đơn SS1 ghi quãng đường mô phỏng theo
thời lượng nhân 15~km/h.
\item \textbf{SS2 và SS5.} Giữ chỗ đỗ đồng thời ghi nhận nhu cầu sạc. Check-in
chuyển chỗ giữ thành vị trí xe đang đỗ, sau đó scheduler cấp cổng nếu còn khả năng.
Nếu chưa cấp được cổng, xe chờ. Check-out đọc thời gian đỗ từ đồng hồ demo và
điện năng, phí quá giờ từ phiên sạc. Hóa đơn được quyết toán một lần trước khi
giải phóng lượt đỗ. SS3 hỗ trợ xác nhận đã đối soát hóa đơn.
\item \textbf{Chính sách mô phỏng.} Giá đỗ mặc định 5.000~đ cho mỗi giờ bắt đầu,
tối thiểu một giờ, cấu hình qua \path{EMOBILITY_PARKING_HOURLY_VND}. Đây là giả định
dùng khi trình diễn. Điện tính 7.000~đ/kWh trong các khung 06:00--09:00 và
16:00--19:30, 3.000~đ/kWh ở giờ khác. Phí chiếm dụng là 10.000~đ mỗi khối
15 phút bắt đầu sau 15 phút ân hạn.
\item \textbf{SS6 và SS7.} Dự báo 24 giờ dùng mô hình 8 lượt phục vụ/giờ và
đơn giá giả định do người vận hành nhập. Preview được tính trên bản sao.
Apply yêu cầu phiên bản live khớp với phiên bản dùng khi tạo preview.
\end{itemize}

\section{Phương thức xử lý bản tin và phiên sạc}
\begin{description}[style=nextline,leftmargin=0pt,labelindent=0pt]
\item[\texttt{feed\_path()} và \texttt{connect\_feed(path)}]
Lấy đường dẫn từ biến môi trường, mở SQLite và tạo bảng sự kiện nếu thiếu.
\item[\texttt{publish\_event(event, path)}]
Kiểm tra trường bắt buộc, loại bản tin và múi giờ của thời điểm phát.
Thêm thời điểm nhận, ghi theo khóa duy nhất \texttt{event\_id}
và trả về trạng thái đã ghi hoặc trùng.
\item[\texttt{consume\_events(state, path)}]
Đọc theo cursor của phiên. Xác minh Hub, chỉ số cổng và thứ tự thời gian.
Ghi kết quả \texttt{APPLIED}, \texttt{STALE} hoặc \texttt{REJECTED}.
\item[\texttt{apply\_hub\_snapshot(state, hubs)}]
Áp dụng snapshot hợp lệ, tăng phiên bản, xử lý phiên sạc bị ảnh hưởng bởi
giảm năng lực và giữ số liệu đã tích lũy.
\item[\texttt{allocate\_hub\_ports(state, hub\_id)}]
Chọn yêu cầu chờ theo điểm ưu tiên, số cổng khả dụng và trần công suất.
Đồng bộ bản ghi xe, lượt đỗ và mã phiên sạc.
\item[\texttt{release\_charging\_session(state, session\_id)}]
Kết thúc phiên đúng mã, tích lũy điện năng và phí sang lượt đỗ nếu có.
Hoàn cổng về số cổng khả dụng trong giới hạn năng lực đang hoạt động.
\end{description}

\section{Giới hạn và bài học}
Trạng thái nghiệp vụ riêng theo phiên giới hạn khả năng dùng chung booking, ví và
lịch sạc giữa nhiều người dùng. Feed chỉ chia sẻ bản tin ngoại vi. Seed lưu số cổng
tổng hợp nên chỉ số trong bản tin chưa ánh xạ chính xác đến ổ cắm của từng phiên sạc.
Đăng nhập, thanh toán, đồng hồ, cảm biến và barie được mô phỏng trong MVP.
Giá đỗ, chất lượng dịch vụ dưới tải lớn và truy cập từ máy ngoài cần xác nhận riêng.
\par
Luồng SS2--SS5 cho thấy điện năng và thời gian chiếm cổng nên được lấy từ phiên sạc,
vì nhập lại tại check-out tạo hai nguồn số liệu cho một hóa đơn. Luồng SS7 cho thấy
preview cần mang phiên bản dữ liệu để thao tác Apply có thể từ chối kết quả đã cũ.
Tách feed IoT khỏi quy trình bấm nút giúp kiểm tra hai cách cập nhật trong cùng demo.

\part{Trình diễn và bàn giao}
\chapter{Trình diễn MVP theo chuỗi màn hình}
\label{chap:demo}
\begin{landscape}
\demoscreen{SS1: Giữ xe dùng chung}
  {Chọn Hub và xe có đủ pin, sau đó xác nhận đặt. Ví ghi khoản cọc 50.000~đ và hệ thống tạo lượt giữ xe.}
  {01-ss1-book.png}{Lượt đặt xe cùng trạng thái giữ cọc.}
\demoscreen{SS1: Nhận, trả và xem hóa đơn}
  {Xác thực PIN trong hạn, chọn Hub trả và hoàn tất quyết toán. Hóa đơn đã lưu hiển thị ở tab trả xe.}
  {02-ss1-invoice.png}{Hóa đơn chuyến đi trong phiên.}
\demoscreen{SS2: Đăng ký, đặt chỗ và check-in}
  {Đăng ký xe cá nhân, đặt chỗ có nhu cầu sạc và xác thực check-in. Trước bước sạc, tab check-out đọc đồng hồ năng lượng từ SS5.}
  {03-ss2-meter.png}{Xe đã vào Hub và điện năng đo từ SS5.}
\demoscreen{SS5: Tiến hành phiên sạc}
  {Mở tab Phiên sạc và Công suất, thực hiện một bước sạc 15 phút. Bảng phiên sạc hiển thị pin, điện năng và trạng thái.}
  {04-ss5-tick.png}{Danh sách phiên sạc sau bước mô phỏng.}
\demoscreen{SS2: Quyết toán và rời trạm}
  {Quay lại SS2, chọn quyết toán và rời trạm. Hệ thống chốt thời gian đỗ, điện năng và phí phát sinh thành một hóa đơn.}
  {05-ss2-invoice.png}{Hóa đơn tổng hợp của xe cá nhân.}
\demoscreen{SS4: Báo cáo sự cố phương tiện}
  {Ghi nhận sự cố từ màn hình Báo cáo và Xử lý sự cố. Bản ghi có mã và trạng thái \texttt{REPORTED}.}
  {06-ss4-incident.png}{Sự cố mới được ghi trong danh sách xử lý.}
\demoscreen{SS6: Xem trước đợt tăng nhu cầu}
  {Chạy Metro Surge, xem sức chứa Hub và các chỉ số trước/sau mô phỏng. Apply là thao tác riêng sau khi xem trước.}
  {07-ss6-preview.png}{So sánh tải của Hub trong preview Metro Surge.}
\demoscreen{SS7: Tạo preview trước bản tin}
  {Tạo Failure Preview và giữ màn hình này. Preview lưu phiên bản trạng thái tại thời điểm chạy mô phỏng.}
  {08-ss7-preview.png}{Bản xem trước sự cố hạ tầng và đề xuất điều hướng.}
\demoscreen{SS3: Cửa sổ thứ nhất nhận cảnh báo}
  {Trong khi SS3 đang mở, tiến trình độc lập gửi sự kiện OFFLINE cho HUB-001, cổng 1. Cảnh báo và mã sự kiện xuất hiện tự động. Lần chạy kiểm chứng ghi mã \path{7985cc2c8f1f4a9eb41eb3e09d5d44d3}.}
  {09-ss3-observer-1.png}{Cửa sổ SS3 thứ nhất quan sát cổng OFFLINE.}
\demoscreen{SS3: Cửa sổ thứ hai nhận cảnh báo}
  {Cửa sổ thứ hai hiển thị cùng mã sự kiện ở trên mà không bấm nút cập nhật. Khi đối chiếu feed trên seed, cổng offline tăng từ 0 lên 1 trong tổng số 10 cổng của HUB-001. Số cổng trống vẫn là 4.}
  {09-ss3-observer-2.png}{Cửa sổ SS3 thứ hai quan sát cùng thay đổi.}
\demoscreen{SS7: Thử áp dụng preview cũ}
  {Quay lại SS7 và bấm Áp dụng. Ứng dụng yêu cầu chạy lại mô phỏng vì phiên bản dữ liệu đã đổi.}
  {10-ss7-stale.png}{Thông báo từ chối và dòng bản tin đi kèm.}
\end{landscape}

\chapter{Khai báo sử dụng AI tạo sinh}
Nhóm BEE BRAVE chịu trách nhiệm về thiết kế, hiện thực và nội dung nộp của Smart E-Mobility Hub. Báo cáo cuối tổng hợp đặc tả và thiết kế S1--S3 của bảy phân hệ, với mỗi thành viên phụ trách một phân hệ. Nhóm dùng ChatGPT và Google Gemini để tham khảo và rà soát kịch bản use case. Codex hỗ trợ biên tập và soạn bản LaTeX từ tài liệu của nhóm, rà soát và chỉnh sửa một số phần mã nguồn, chạy kiểm thử và chụp ảnh demo. Nhóm quyết định nội dung sử dụng trong bản nộp.

\end{document}
